\textbf{Phần 1. Trắc nghiệm lựa chọn}
\setcounterpageref{ex}{0}

\begin{ex} %[5D1N5]
Kết quả của phép tính $\dfrac{6}{5} + \dfrac{7}{15}$ là
\choice
{$\dfrac{13}{20}$}
{$\dfrac{13}{15}$}
{$\dfrac{25}{30}$}
{\True $\dfrac{5}{3}$}
\loigiai{
Quy đồng mẫu số: $\dfrac{6}{5} + \dfrac{7}{15} = \dfrac{18}{15} + \dfrac{7}{15} = \dfrac{25}{15} = \dfrac{5}{3}$.\\
\textbf{Đáp án của câu hỏi là D}
}
\end{ex}

\begin{ex} %[5D1N5]
Kết quả của phép trừ $\dfrac{15}{16} - \dfrac{3}{4}$ là
\choice
{\True $\dfrac{3}{16}$}
{$\dfrac{12}{12}$}
{$\dfrac{12}{16}$}
{$\dfrac{9}{16}$}
\loigiai{
Quy đồng mẫu số: $\dfrac{15}{16} - \dfrac{3}{4} = \dfrac{15}{16} - \dfrac{12}{16} = \dfrac{3}{16}$.\\
\textbf{Đáp án của câu hỏi là A}
}
\end{ex}

\begin{ex} %[5D1H5]
Kết quả của phép tính $3 - \dfrac{16}{7}$ là
\choice
{$\dfrac{13}{7}$}
{$\dfrac{21}{7}$}
{\True $\dfrac{5}{7}$}
{$\dfrac{1}{7}$}
\loigiai{
Ta có: $3 - \dfrac{16}{7} = \dfrac{21}{7} - \dfrac{16}{7} = \dfrac{5}{7}$.\\
\textbf{Đáp án của câu hỏi là C}
}
\end{ex}

\begin{ex} %[5D1H5]
Kết quả của phép tính $\dfrac{5}{8} \times \dfrac{8}{15}$ là
\choice
{$\dfrac{1}{5}$}
{\True $\dfrac{1}{3}$}
{$\dfrac{5}{15}$}
{$\dfrac{3}{8}$}
\loigiai{
Rút gọn tích: $\dfrac{5 \times 8}{8 \times 15} = \dfrac{5}{15} = \dfrac{1}{3}$.\\
\textbf{Đáp án của câu hỏi là B}
}
\end{ex}

\begin{ex} %[5D1H5]
Kết quả của phép chia $\dfrac{12}{7} : \dfrac{8}{7}$ là
\choice
{$\dfrac{4}{7}$}
{$\dfrac{2}{3}$}
{\True $\dfrac{3}{2}$}
{$\dfrac{96}{49}$}
\loigiai{
Thực hiện phép chia: $\dfrac{12}{7} : \dfrac{8}{7} = \dfrac{12}{7} \times \dfrac{7}{8} = \dfrac{12}{8} = \dfrac{3}{2}$.\\
\textbf{Đáp án của câu hỏi là C}
}
\end{ex}

\begin{ex} %[5D1H5]
Kết quả của phép tính $\dfrac{24}{5} : 4$ là
\choice
{$\dfrac{96}{5}$}
{$\dfrac{24}{20}$}
{$\dfrac{1}{5}$}
{\True $\dfrac{6}{5}$}
\loigiai{
Ta có: $\dfrac{24}{5} : 4 = \dfrac{24 : 4}{5} = \dfrac{6}{5}$.\\
\textbf{Đáp án của câu hỏi là D}
}
\end{ex}

\begin{ex} %[5D1H5]
Giá trị của biểu thức $\dfrac{4}{7} + \dfrac{9}{14} - \dfrac{5}{7}$ là
\choice
{\True $\dfrac{1}{2}$}
{$\dfrac{1}{14}$}
{$\dfrac{8}{14}$}
{$\dfrac{3}{7}$}
\loigiai{
Giao hoán và kết hợp: $\left(\dfrac{4}{7} - \dfrac{5}{7}\right) + \dfrac{9}{14} = \dfrac{9}{14} - \dfrac{1}{7} = \dfrac{9}{14} - \dfrac{2}{14} = \dfrac{7}{14} = \dfrac{1}{2}$.\\
\textbf{Đáp án của câu hỏi là A}
}
\end{ex}

\begin{ex} %[5D1H5]
Giá trị của biểu thức $\dfrac{9}{10} \times \dfrac{5}{6} : 3$ là
\choice
{$\dfrac{3}{4}$}
{$\dfrac{9}{4}$}
{\True $\dfrac{1}{4}$}
{$\dfrac{1}{2}$}
\loigiai{
Thực hiện lần lượt từ trái sang phải: $\dfrac{9 \times 5}{10 \times 6} = \dfrac{45}{60} = \dfrac{3}{4}$.\\
Sau đó: $\dfrac{3}{4} : 3 = \dfrac{1}{4}$.\\
\textbf{Đáp án của câu hỏi là C}
}
\end{ex}

\begin{ex} %[5D1H5]
Tính nhanh biểu thức $\dfrac{10}{11} \times \dfrac{9}{16} \times \dfrac{11}{10} \times \dfrac{8}{9}$ ta được kết quả là
\choice
{$1$}
{\True $\dfrac{1}{2}$}
{$2$}
{$\dfrac{1}{4}$}
\loigiai{
Nhóm các phân số nghịch đảo và rút gọn:
\[
\left(\dfrac{10}{11} \times \dfrac{11}{10}\right) \times \left(\dfrac{9}{16} \times \dfrac{8}{9}\right) = 1 \times \dfrac{8}{16} = \dfrac{1}{2}.
\]
\textbf{Đáp án của câu hỏi là B}
}
\end{ex}

\begin{ex} %[5D1V5]
Hai bạn Nam và Việt cùng chạy thi trên một đoạn đường. Sau một phút, Nam chạy được $\dfrac{3}{5}$ đoạn đường, Việt chạy được $\dfrac{7}{10}$ đoạn đường. Khẳng định nào sau đây là đúng?
\choice
{Nam chạy nhiều hơn Việt $\dfrac{1}{10}$ đoạn đường}
{Việt chạy nhiều hơn Nam $\dfrac{4}{10}$ đoạn đường}
{Hai bạn chạy được quãng đường bằng nhau}
{\True Việt chạy nhiều hơn Nam $\dfrac{1}{10}$ đoạn đường}
\loigiai{
Quy đồng: $\dfrac{3}{5} = \dfrac{6}{10}$.\\
Vì $\dfrac{7}{10} > \dfrac{6}{10}$ nên Việt chạy được nhiều hơn Nam là: $\dfrac{7}{10} - \dfrac{6}{10} = \dfrac{1}{10}$ đoạn đường.\\
\textbf{Đáp án của câu hỏi là D}
}
\end{ex}

\begin{ex} %[5D1V5]
Thư viện trường học có $\dfrac{2}{3}$ số sách là sách giáo khoa, $\dfrac{2}{9}$ số sách là sách tham khảo, còn lại là truyện thiếu nhi và tạp chí. Phân số chỉ số truyện thiếu nhi và tạp chí là
\choice
{\True $\dfrac{1}{9}$}
{$\dfrac{4}{9}$}
{$\dfrac{8}{9}$}
{$\dfrac{2}{9}$}
\loigiai{
Tổng số phần sách giáo khoa và sách tham khảo là: $\dfrac{2}{3} + \dfrac{2}{9} = \dfrac{6}{9} + \dfrac{2}{9} = \dfrac{8}{9}$.\\
Phân số chỉ số truyện thiếu nhi và tạp chí là: $1 - \dfrac{8}{9} = \dfrac{1}{9}$.\\
\textbf{Đáp án của câu hỏi là A}
}
\end{ex}

\begin{ex} %[5D1V5]
Một tấm kính hình chữ nhật có chiều dài $\dfrac{5}{2}\text{ m}$, chiều rộng $\dfrac{4}{3}\text{ m}$. Chú Hoà chia tấm kính đó thành $3$ phần bằng nhau để làm mặt bàn. Diện tích mỗi phần tấm kính làm mặt bàn là
\choice
{$\dfrac{10}{3}\text{ m}^2$}
{\True $\dfrac{10}{9}\text{ m}^2$}
{$\dfrac{5}{9}\text{ m}^2$}
{$\dfrac{20}{9}\text{ m}^2$}
\loigiai{
Diện tích cả tấm kính là: $\dfrac{5}{2} \times \dfrac{4}{3} = \dfrac{20}{6} = \dfrac{10}{3}\text{ (m}^2\text{)}$.\\
Diện tích mỗi phần làm mặt bàn là: $\dfrac{10}{3} : 3 = \dfrac{10}{9}\text{ (m}^2\text{)}$.\\
\textbf{Đáp án của câu hỏi là B}
}
\end{ex}

\textbf{Phần 2. Trắc nghiệm đúng sai}
\setcounterpageref{ex}{0}

\begin{ex} %[5D1H5]
Cho các biểu thức: $A = \dfrac{29}{12} - \left(2 + \dfrac{1}{4}\right)$, $B = \dfrac{5}{6} \times \dfrac{3}{7} + \dfrac{3}{7} \times \dfrac{7}{6}$ và $C = \dfrac{11}{4} : \left(\dfrac{11}{14} \times 7\right)$. Xét tính đúng sai của các khẳng định sau:
\choiceTF
{\True Biểu thức $A$ có giá trị bằng $\dfrac{1}{6}$.}
{Biểu thức $B$ có giá trị bằng $1$.}
{\True Biểu thức $C$ có giá trị bằng $\dfrac{1}{2}$.}
{\True Giá trị của biểu thức $B$ lớn hơn giá trị của biểu thức $C$.}
\loigiai{
\begin{itemize}
    \item $A = \dfrac{29}{12} - \dfrac{9}{4} = \dfrac{29}{12} - \dfrac{27}{12} = \dfrac{2}{12} = \dfrac{1}{6}$.
    \item $B = \dfrac{3}{7} \times \left(\dfrac{5}{6} + \dfrac{7}{6}\right) = \dfrac{3}{7} \times 2 = \dfrac{6}{7} \neq 1$.
    \item $C = \dfrac{11}{4} : \left(\dfrac{11}{2}\right) = \dfrac{11}{4} \times \dfrac{2}{11} = \dfrac{2}{4} = \dfrac{1}{2}$.
    \item So sánh $B = \dfrac{6}{7} \approx 0{,}86$ và $C = \dfrac{1}{2} = 0{,}5 \Rightarrow B > C$.
\end{itemize}
\begin{itemize}
    \item a) \textbf{Đúng}.
    \item b) \textbf{Sai}: $B = \dfrac{6}{7}$.
    \item c) \textbf{Đúng}.
    \item d) \textbf{Đúng}.
\end{itemize}
}
\end{ex}

\begin{ex} %[5D1V5]
Hưởng ứng phong trào quyên góp sách vào thư viện trường học, lớp 5A quyên góp được $96$ quyển sách, lớp 5B quyên góp được số quyển sách bằng $\dfrac{7}{8}$ số quyển sách của lớp 5A. Xét tính đúng sai của các khẳng định sau:
\choiceTF
{\True Lớp 5B quyên góp được $84$ quyển sách.}
{Lớp 5B quyên góp được nhiều sách hơn lớp 5A.}
{\True Cả hai lớp 5A và 5B quyên góp được tất cả $180$ quyển sách.}
{Lớp 5A quyên góp nhiều hơn lớp 5B là $14$ quyển sách.}
\loigiai{
\begin{itemize}
    \item Số sách lớp 5B quyên góp là: $96 \times \dfrac{7}{8} = 84$ (quyển).
    \item Vì $84 < 96$ nên lớp 5B quyên góp ít hơn lớp 5A.
    \item Cả hai lớp quyên góp: $96 + 84 = 180$ (quyển).
    \item Lớp 5A quyên góp nhiều hơn lớp 5B là: $96 - 84 = 12$ quyển ($\neq 14$ quyển).
\end{itemize}
\begin{itemize}
    \item a) \textbf{Đúng}.
    \item b) \textbf{Sai}: Lớp 5B quyên góp ít hơn lớp 5A.
    \item c) \textbf{Đúng}.
    \item d) \textbf{Sai}: Nhiều hơn $12$ quyển sách.
\end{itemize}
}
\end{ex}

\textbf{Phần 3. Trả lời ngắn}
\setcounterpageref{ex}{0}

\begin{ex} %[5D1H5]
Tính giá trị của biểu thức sau dưới dạng phân số tối giản $\dfrac{a}{b}$:
\[
A = \dfrac{5}{6} \times \dfrac{3}{7} + \dfrac{3}{7} \times \dfrac{7}{6}.
\]
Tính giá trị của $a + b$.
\shortans{13}
\loigiai{
\textbf{Đáp số:} $13$\\
$A = \dfrac{3}{7} \times \left(\dfrac{5}{6} + \dfrac{7}{6}\right) = \dfrac{3}{7} \times 2 = \dfrac{6}{7}$. Do đó $a = 6, b = 7 \Rightarrow a + b = 13$.
}
\end{ex}

\begin{ex} %[5D1H5]
Người ta gắn dây đèn xung quanh một tấm biển quảng cáo hình vuông có chu vi là $18\text{ m}$. Tính diện tích của tấm biển quảng cáo đó dưới dạng phân số tối giản $\dfrac{m}{n}\text{ (m}^2\text{)}$. Giá trị của tử số $m$ là bao nhiêu?
\shortans{81}
\loigiai{
\textbf{Đáp số:} $81$\\
Cạnh của hình vuông là: $18 : 4 = \dfrac{9}{2}\text{ (m)}$.\\
Diện tích tấm biển là: $\dfrac{9}{2} \times \dfrac{9}{2} = \dfrac{81}{4}\text{ (m}^2\text{)}$. Vậy $m = 81$.
}
\end{ex}

\begin{ex} %[5D1V5]
Lớp 5A quyên góp được $96$ quyển sách, lớp 5B quyên góp được số quyển sách bằng $\dfrac{7}{8}$ số sách của lớp 5A. Hỏi cả hai lớp đã quyên góp được tất cả bao nhiêu quyển sách?
\shortans{180}
\loigiai{
\textbf{Đáp số:} $180$\\
Số sách lớp 5B quyên góp là: $96 \times \dfrac{7}{8} = 84$ quyển.\\
Cả hai lớp quyên góp: $96 + 84 = 180$ quyển.
}
\end{ex}

\begin{ex} %[5D1V5]
Tấm kính hình chữ nhật có chiều dài $\dfrac{5}{2}\text{ m}$, chiều rộng $\dfrac{4}{3}\text{ m}$ được chia đều làm $3$ mặt bàn bằng nhau. Diện tích của mỗi mặt bàn là $\dfrac{p}{q}\text{ m}^2$ (với $\dfrac{p}{q}$ là phân số tối giản). Tính tổng $p + q$.
\shortans{19}
\loigiai{
\textbf{Đáp số:} $19$\\
Diện tích cả tấm kính: $\dfrac{5}{2} \times \dfrac{4}{3} = \dfrac{10}{3}\text{ m}^2$.\\
Diện tích mỗi mặt bàn: $\dfrac{10}{3} : 3 = \dfrac{10}{9}\text{ m}^2 \Rightarrow p = 10, q = 9 \Rightarrow p + q = 19$.
}
\end{ex}

\begin{ex} %[5D1H5]
Tính giá trị của biểu thức:
\[
M = \dfrac{9}{7} \times \left(\dfrac{7}{9} - \dfrac{2}{3}\right).
\]
Viết kết quả dưới dạng phân số tối giản $\dfrac{1}{k}$. Tìm giá trị của $k$.
\shortans{7}
\loigiai{
\textbf{Đáp số:} $7$\\
$M = \dfrac{9}{7} \times \left(\dfrac{7}{9} - \dfrac{6}{9}\right) = \dfrac{9}{7} \times \dfrac{1}{9} = \dfrac{1}{7} \Rightarrow k = 7$.
}
\end{ex}

\begin{ex} %[5D1V5]
Để cắt lấy đoạn băng giấy dài $\dfrac{1}{2}\text{ m}$ từ một băng giấy ban đầu dài $\dfrac{2}{3}\text{ m}$ mà không dùng thước đo, ta cần cắt bỏ đi một đoạn băng giấy có chiều dài bằng bao nhiêu mét? (Viết kết quả dưới dạng phân số tối giản $\dfrac{1}{x}$). Tìm giá trị của $x$.
\shortans{6}
\loigiai{
\textbf{Đáp số:} $6$\\
Đoạn băng giấy cần cắt bỏ có chiều dài là: $\dfrac{2}{3} - \dfrac{1}{2} = \dfrac{4}{6} - \dfrac{3}{6} = \dfrac{1}{6}\text{ (m)}$. Do đó $x = 6$.
}
\end{ex}
